%File: formatting-instructions-latex-2024.tex
%release 2024.0
\documentclass[letterpaper]{article} % DO NOT CHANGE THIS
\usepackage{aaai24}  % DO NOT CHANGE THIS
\usepackage{times}  % DO NOT CHANGE THIS
\usepackage{helvet}  % DO NOT CHANGE THIS
\usepackage{courier}  % DO NOT CHANGE THIS
\usepackage[hyphens]{url}  % DO NOT CHANGE THIS
\usepackage{graphicx} % DO NOT CHANGE THIS
\urlstyle{rm} % DO NOT CHANGE THIS
\def\UrlFont{\rm}  % DO NOT CHANGE THIS
\usepackage{natbib}  % DO NOT CHANGE THIS AND DO NOT ADD ANY OPTIONS TO IT
\usepackage{caption} % DO NOT CHANGE THIS AND DO NOT ADD ANY OPTIONS TO IT
\frenchspacing  % DO NOT CHANGE THIS
\setlength{\pdfpagewidth}{8.5in}  % DO NOT CHANGE THIS
\setlength{\pdfpageheight}{11in}  % DO NOT CHANGE THIS
%
% These are recommended to typeset algorithms but not required. See the subsubsection on algorithms. Remove them if you don't have algorithms in your paper.
\usepackage{algorithm}
\usepackage{algorithmic}
\usepackage{booktabs}  % For professional-quality tables
%
% These are are recommended to typeset listings but not required. See the subsubsection on listing. Remove this block if you don't have listings in your paper.
\usepackage{newfloat}
\usepackage{listings}
\DeclareCaptionStyle{ruled}{labelfont=normalfont,labelsep=colon,strut=off} % DO NOT CHANGE THIS
\lstset{%
	basicstyle={\footnotesize\ttfamily},% footnotesize acceptable for monospace
	numbers=left,numberstyle=\footnotesize,xleftmargin=2em,% show line numbers, remove this entire line if you don't want the numbers.
	aboveskip=0pt,belowskip=0pt,%
	showstringspaces=false,tabsize=2,breaklines=true}
\floatstyle{ruled}
\newfloat{listing}{tb}{lst}{}
\floatname{listing}{Listing}
%
% Keep the \pdfinfo as shown here. There's no need
% for you to add the /Title and /Author tags.
\pdfinfo{
/TemplateVersion (2024.1)
}

% DISALLOWED PACKAGES
% \usepackage{authblk} -- This package is specifically forbidden
% \usepackage{balance} -- This package is specifically forbidden
% \usepackage{color (if used in text)
% \usepackage{CJK} -- This package is specifically forbidden
% \usepackage{float} -- This package is specifically forbidden
% \usepackage{flushend} -- This package is specifically forbidden
% \usepackage{fontenc} -- This package is specifically forbidden
% \usepackage{fullpage} -- This package is specifically forbidden
% \usepackage{geometry} -- This package is specifically forbidden
% \usepackage{grffile} -- This package is specifically forbidden
% \usepackage{hyperref} -- This package is specifically forbidden
% \usepackage{navigator} -- This package is specifically forbidden
% (or any other package that embeds links such as navigator or hyperref)
% \indentfirst} -- This package is specifically forbidden
% \layout} -- This package is specifically forbidden
% \multicol} -- This package is specifically forbidden
% \nameref} -- This package is specifically forbidden
% \usepackage{savetrees} -- This package is specifically forbidden
% \usepackage{setspace} -- This package is specifically forbidden
% \usepackage{stfloats} -- This package is specifically forbidden
% \usepackage{tabu} -- This package is specifically forbidden
% \usepackage{titlesec} -- This package is specifically forbidden
% \usepackage{tocbibind} -- This package is specifically forbidden
% \usepackage{ulem} -- This package is specifically forbidden
% \usepackage{wrapfig} -- This package is specifically forbidden
% DISALLOWED COMMANDS
% \nocopyright -- Your paper will not be published if you use this command
% \addtolength -- This command may not be used
% \balance -- This command may not be used
% \baselinestretch -- Your paper will not be published if you use this command
% \clearpage -- No page breaks of any kind may be used for the final version of your paper
% \columnsep -- This command may not be used
% \newpage -- No page breaks of any kind may be used for the final version of your paper
% \pagebreak -- No page breaks of any kind may be used for the final version of your paperr
% \pagestyle -- This command may not be used
% \tiny -- This is not an acceptable font size.
% \vspace{- -- No negative value may be used in proximity of a caption, figure, table, section, subsection, subsubsection, or reference
% \vskip{- -- No negative value may be used to alter spacing above or below a caption, figure, table, section, subsection, subsubsection, or reference

\setcounter{secnumdepth}{0} %May be changed to 1 or 2 if section numbers are desired.

% The file aaai24.sty is the style file for AAAI Press
% proceedings, working notes, and technical reports.
%

\title{Machine Learning for Exoplanet Detection: Identifying Exoplanets in Light Curves}
\author {
    % Authors
    Tristan Moffett\textsuperscript{\rm },
    Josh Manchester\textsuperscript{\rm },
    Brianne Leatherman\textsuperscript{\rm }
}
\affiliations {
    % Affiliations
    \textsuperscript{\rm }University of Colorado at Colorado Springs\\
    %\textsuperscript{\rm 2}Affiliation 2\\
    %firstAuthor@affiliation1.com, secondAuthor@affilation2.com, thirdAuthor@affiliation1.com
}

\begin{document}
\maketitle

\section*{Abstract}
The discovery of exoplanets has become more dependent on machine learning as missions like Kepler and TESS generate thousands of transit signals that are too costly and time consuming to vet by hand. In this research, we address this problem by comparing three different deep learning models for exoplanet classification: Convolutional Neural Networks (CNNs), Recurrent Neural Networks (RNNs), and Transformer architecture. Each team member designed, trained, and evaluated one of these models using observational data from Kepler and the Transiting Exoplanet Survey Satellite (TESS) along with simulated datasets. Our focus is on performance metrics such as accuracy, precision, recall, F1 score, and Area Under Curve (AUC), as well as understanding the strengths and weaknesses of each model in classifying planets versus false positives.

This paper presents preliminary results from the BiLSTM + clustering approach, which achieved an AUC of 0.6947 on validation data (655 windows with 150 positive examples) and successfully identified TIC 307210830, a confirmed multi-planet system from TESS, demonstrating real-world applicability. The RNN implementation combines K-means clustering on Box Least Squares features with a 3-layer bidirectional LSTM architecture, enabling the network to learn cluster-specific patterns and distinguish different types of stellar variability from genuine planetary transits. The model uses class weighting to handle the 23 percent positive class imbalance and achieves meaningful predictive power despite the small training dataset. From analyzing results across the three architectures, we aim to identify which methods are most effective for generalizability across multiple datasets and advance automated exoplanet vetting for future missions.



\section*{Introduction \& Background}
In the last twenty years, there have been significant increases in the number of space missions whose goal is to observe distant stars and the changes that may occur over variable lengths of time. NASA has been one of the primary drivers behind these advances, which are called transit survey missions. Of most importance to our research are two missions in particular: Kepler/K2 and TESS. Kepler was launched in 2009 with the original mission of monitoring a 116 square degree segment of sky for three and a half years, and when the satellite had a mechanical issue, the mission was changed to K2, which monitored stars and systems along the elliptic plane until it was retired in 2018 \cite{salinas2023distinguishingtransitfalsepositives}. The Transiting Exoplanet Survey Satellite (TESS) was launched in 2018, nearly ten years after Kepler, with the goal of monitoring a sector of approximately 2300 square degrees for about 27 days before moving to the next section \cite{salinas2023distinguishingtransitfalsepositives}. At the time of this writing, these two transit surveys have accounted for more than 3000 of the 6000 total confirmed exoplanets \cite{salinas2025tessfullframe}.

The most prevalent method of determining the presence of an extra-orbital planet has been through observing changes in the intensity of the brightness of a star when an object transits between the star and the observer. These transits can be found by representing the observed flux from the star in a time-series graph, also known as light curves, and then analyzing the data \cite{morvan2022dontpayattentionnoise}. This process often includes phase-folding the data by overlapping periods to determine patterns \cite{salinas2023distinguishingtransitfalsepositives}. A dip in the light curve indicates a threshold crossing event (TCE) but that alone does not mean a planet caused the obstruction of the light. The TCE can also be a false positive that was triggered by an eclipsing binary or instrumental artifact \cite{salinas2025tessfullframe}.

The similar nature of the light curve of a transiting exoplanet and that of an eclipsing binary can be difficult to differentiate. Additionally, light curves can be challenging to work with due to the continuous nature of background noise from instrumental, photon, and other interference in the transit of light, and the curves can require significant work to pre-process and then analyze the data \cite{morvan2022dontpayattentionnoise}. As deep learning continues to develop, the use of neural networks has become a common tool in this process, significantly reducing the number of light curves that need to be analyzed by hand.

The Convolutional Neural Network (CNN) was introduced as a model designed to detect repeating patterns in data. For exoplanet research, this means recognizing the U-shaped dips of planetary transits and separating them from the V-shaped dips of eclipsing binaries within time-series light curve data. A challenge for CNNs is that light curves often contain a large amount of noise, which can make training more difficult. However, when handled properly, CNNs can still learn to use this noise to improve classification. \cite{Dattilo2019identifyexoplanetsII} demonstrated the value of CNNs on Kepler data by combining global and local views of a light curve to capture both the orbital view and the details of a transit. \cite{Osborn2019rapidclassificationtesscandidates} later applied a similar approach to TESS data and introduced additional features such as centroid and stellar parameters to handle noisier signals. These results show why CNNs are a strong fit for exoplanet classification, though other deep learning approaches like Recurrent Neural Networks (RNNs) and Transformers bring their own strengths for working with time-series data such as light curves.

Recurrent Neural Networks (RNNs) are a natural fit for time-series data because they learn patterns across many time steps and keep context memory (history) from earlier points in the sequence. In light-curve analysis, this lets the model use both the local transit shape of the light curve and the surrounding temporal pattern (periodic repeats) at the same time. Gated variants like LSTM and GRU help with long-range dependencies and noisy sequences of data. Stacked LSTMs have shown to be a strong resource for space-mission photometry: for example, an LSTM detector trained on Kepler flares generalized well to TESS with high precision/recall when preprocessing and class balancing were handled carefully \cite{Vida2021findingflaresrecurrent}. Beyond direct classification, sequence-aware representation learning with an Echo State Network (ESN) + autoencoder reconstructs light curves in sequence space (not just a readout vector), producing more faithful time-series embeddings for future tasks \cite{kugler2016explorativeapproachkepler}. Finally, RMTPP shows how an RNN can encode event history to predict both the timing and type of the next event; adding a lightweight timing cue (intervals/slight inaccurate period) alongside the classifier can reduce aperiodic false positives in transit searches \cite{Du2016markedtemporalpoint}. In this project, we will benchmark LSTM/GRU on our curated subset, report precision, recall, F1, and ROC-AUC, and optionally include a small timing helper to capture periodicity.

The Transformer architecture was introduced as a novel, simple network relying on encoder/decoder layers of self-attention mechanisms \cite{vaswani2017attentionneed}. Because attention can be calculated in parallel, this network saw an improvement in quality of results with a reduced amount of training time compared to prior RNN and CNN approaches \cite{vaswani2017attentionneed}. Although Transformers were originally designed for Natural Language Processing (NLP), their use in other fields has rapidly spread, due to both the efficiency and the ability to correlate sequential data like time series such as a light curve, where the transit of a planet can be considered a function of time \cite{salinas2023distinguishingtransitfalsepositives}. Given the rapid advancements in deep learning methods over the last decade, a direct performance comparison across RNNs, CNNs, and Transformers on light curve analysis should prove interesting. This paper will detail the methods for creating such a comparison, as well as results, by modeling multiple research processes from other papers.


\section*{Related Work}
\subsection{Convolutional Neural Networks (CNNs)}
Shallue and Vanderburg (\citeyear{Shallue2018identifyexoplanetsI}) addressed the problem of classifying Threshold Crossing Events (TCEs) from the Kepler mission, where thousands of signals had to be examined and most turned out not to be planets. To resolve this bottleneck, they designed one of the first convolutional neural networks (CNNs) for exoplanet detection, comparing a linear model, a fully connected network, and a 1-D CNN, with the CNN proving most successful. Their final architecture called AstroNet processed two versions of each light curve: a global view covering the full orbital phase and a local view zoomed in around the transit. Each view passed through layers of one-dimensional convolutions and max pooling before being merged into fully connected layers of 512 nodes, with dropout and data augmentation used to reduce overfitting. On Kepler's real observational test set the CNN achieved results with 96\% accuracy, 98.8\% AUC, 96\% precision, and 95\% recall. However, when evaluated on simulated datasets the model struggled, recovering 81\% of true injected planets but incorrectly labeling around 90\% of background planets and eclipsing binaries as planets, misclassifying about 65\% of eclipsing binaries, and letting through 3.18\% of false positives from inverted light curves. These outcomes demonstrated that while the model was powerful on Kepler data, it lacked generalizability, leading the authors to suggest training with simulated data and including centroid information to better handle false positives. This highlights the need for models that can work across both real and simulated datasets, a direction we plan to take in our own work.

Dattilo et al. (\citeyear{Dattilo2019identifyexoplanetsII}) focused on the challenge of vetting the large number of planet candidates identified in K2 light curves, which are especially noisy due to pointing instabilities from the Kepler telescope. This created difficulty in classifying true planet signals from false positives such as eclipsing binaries or systematics. To address this, they extended the AstroNet CNN first built for Kepler data by Shallue and Vanderburg (2018), adapting it to the shorter and noisier K2 campaigns. Their approach used both real K2 TCEs labeled as candidates, binaries, or junk, and tested on simulated ``Kepler-lite'' data to tune their model under similar conditions. Each TCE was represented by two views of the light curve, a global view showing the full orbital phase and a local view zoomed in on the transit. These were normalized and then passed through separate convolutional and max pooling layers before being merged into fully connected layers of 512 nodes, with additional scalar features such as planet-to-star radius ratio and impact parameter included to strengthen classification. The final output gave the probability of a candidate being a true planet. Their CNN achieved high accuracy averaging 98\% on test sets and successfully uncovered two new super-Earths within the K2 dataset. This work highlights the need for both global and local inputs including additional scalar features as it improves against noisy data while also improving the classification of planets and false positives.

Osborne et al. (\citeyear{Osborn2019rapidclassificationtesscandidates}) addressed the challenge of rapidly classifying the thousands of planet candidates generated monthly by TESS. To further the issue each month around 500 hours were required to select the best candidates from each month's findings. As any sane person would be able to tell, this is a bottleneck as it requires too many resources and people that could be used elsewhere. To resolve this issue however, Osborne and his colleagues trained a CNN model based on Shallue and Vanderburg (2018) and Ansdell et al. (2018) work, specifically the AstroNet and ExoNet CNN models. They trained and validated their CNN strictly on high-fidelity simulated TESS light curves (the TSOP-301 dataset), which provided known ground truths for both planets and false positives. Each candidate was represented by both a local and global view of its light curve, with centroid information also included to help spot background eclipsing binaries. Alongside this, during training, they included balanced batch sampling to counter class imbalance (only 14\% of candidates were true planets) and applied augmentation techniques to improve generalization and reduce overfitting. When tested, their model achieved strong results on their simulated set, with a 97\% average precision in binary classification and a 92\% accuracy in identifying planets in their multi-class classification. However, when using the real TESS set, the model's performance dropped, only being able to recover 61\% TOIs (TESS Objects of Interest) as planets. Though the model was still able to flag around 200 new TCEs as possible planets. Results highlight a reduction in human workload and planet discovery but also show a gap between simulated and real data performance in CNNs. This brings a need to build a model that is more generalizable capable of handling both simulated and observational dataset together, which we plan to explore in our project.

\subsubsection{CNN Classifier for Next Generation Transit Survey}
Chaushev et al. (\citeyear{Chaushev2019ngts}) focused on improving the speed and consistency of exoplanet detection in large transit surveys, particularly the Next Generation Transit Survey (NGTS), where manual vetting of light curves is reliable but extremely time-consuming. To automate this process, the authors introduced a convolutional neural network (CNN) called PlaNET, which aimed to reduce dependence on large labeled datasets by using a mix of real NGTS observations, unlabeled data, and simulated planetary transits. This combination of data helped the model learn from both real and synthetic signals while improving generalization to noisy and incomplete data. Each light curve was detrended, normalized, and phase-folded so that transits aligned, creating a consistent input for classification. Their architecture was inspired by AstroNet, featuring both global and local views of the light curve to capture the full orbital shape and transit depths. Additional stellar information, such as radius and orbital period, was included to help the model separate planets from false positives like eclipsing binaries or background blends. The model was trained using the Adam optimizer with early stopping and 10-fold cross-validation to avoid overfitting. In testing, PlaNET achieved around 96\% accuracy and an AUC of 99.3\%, successfully recovering 13 out of 14 confirmed NGTS planets while identifying several new candidates missed by human vetting. Their results show that CNNs can greatly reduce manual effort and improve the reliability of large-scale exoplanet searches. The main takeaway here is the use of extra features such as the radius and orbital period. These features further help in the process of classification, especially in the case where exoplanet lightcurves look similar to EBs and BEBs.

\subsubsection{Stella CNN Flare Classifier}
Feinstein et al. (\citeyear{Feinstein2020stella}) built a CNN called stella to detect stellar flares in TESS short-cadence data. The problem they faced was that manually identifying flares across thousands of young stars is too time-consuming, especially since high stellar variability makes simple outlier detection methods fail. Their solution was to train a CNN to recognize flare shapes directly from the light curves and use this to study how flare rates change with stellar age and temperature. They worked with 3,200 young stars from TESS Sectors 1–20, detrending and normalizing each light curve before splitting them into 200-cadence segments centered on either known flare events or non-flare regions. The architecture used standard components with one-dimensional convolutional layers to extract features, pooling layers to reduce dimensions, fully connected layers with dropout to prevent overfitting, and a final sigmoid layer for binary classification. They trained it with the Adam optimizer and binary cross-entropy loss, running 10 models with different initializations to create an ensemble and reduce overfitting. Their model stella performed really well, hitting an AUC of 0.993 and F1 scores above 0.98 on validation data. When applied to the full sample, it found over 17,000 flares and showed that flare activity drops off significantly for stars older than 50 Myr and hotter than 4000 K. What's interesting here is the data preparation approach. Instead of phase-folding like most exoplanet detection methods, they used fixed-length segments centered on events of interest. This could also work for planet detection, where we would just center segments on threshold crossing events instead of flares.

\subsubsection{GPU-accelerated Phase Folding}
Wang et al. (\citeyear{Wang2024gpfc}) developed GPFC, a pipeline that combines GPU-accelerated phase folding with a CNN to detect ultra-short-period planets with orbital periods less than one day. Most exoplanet detection pipelines rely on Box Least Squares (BLS) to identify transits, but BLS becomes extremely slow when scanning through thousands of trial periods at high resolution. Wang et al. took a different route by offloading the phase folding computation to the GPU and using a CNN to score each folded result. The workflow starts with detrending and outlier removal, then creates a dense grid of trial periods to test. Next, the light curve gets folded and binned at each period, followed by noise normalization to standardize the input. Finally, a CNN evaluates each folded light curve and outputs a score indicating the likelihood of a real transit signal. The impressive part is the speed at which they can fold a single light curve at 100,000 different trial periods in just five seconds on an RTX 3090 Ti. This works by batching 16-period folds simultaneously and using efficient binning algorithms that take advantage of GPU parallelization. The CNN itself has 19 layers and includes circular convolution, global max pooling, and dropout layers to prevent overfitting. They trained it using TensorFlow with the Adam optimizer. When they tested GPFC against traditional BLS methods, it performed better at detecting low signal-to-noise ratio transits and ran much faster. In a blind search over the 0.2-1.0 day period range, GPFC successfully recovered all confirmed ultra-short-period planets in the Kepler catalog. What makes this approach valuable is that GPU acceleration turns exhaustive period searches from impractical to practical. Testing 100,000 trial periods per light curve would take hours with the usual BLS method, but the GPU implementation cuts this down to seconds. Our current pipeline doesn't use GPU-accelerated folding yet, but this technique could become important if we need to process larger datasets or want to search finer period grids in future work.

\subsubsection{Applying These Ideas to Our Work}
Looking at what worked across these papers, we're going to borrow several techniques that consistently improve detection performance. The dual-view architecture from Shallue and Vanderburg (\citeyear{Shallue2018identifyexoplanetsI}) is one of the most useful ideas. It splits each light curve into a global view (2001 bins) covering the full orbital phase and a local view (201 bins) zoomed in on the transit itself. This way, the CNN gets both the orbital context and the finer transit details. We'll follow their preprocessing approach too: detrend with median subtraction, use Box Least Squares to find the period, phase fold around the transit center, and normalize so the baseline sits at zero with the deepest transit point at -1.

The auxiliary scalar features from Dattilo et al. (\citeyear{Dattilo2019identifyexoplanetsII}) and Chaushev et al. (\citeyear{Chaushev2019ngts}) seem worth including as well. Adding things like planet-to-star radius ratio, impact parameter, transit depth, and orbital period helps the model distinguish real planets from eclipsing binaries, especially when the light curves alone are ambiguous or noisy. Osborne et al. (\citeyear{Osborn2019rapidclassificationtesscandidates}) showed that centroid information can catch background eclipsing binaries that might otherwise slip through, and their balanced batch sampling approach handles the severe class imbalance problem we're likely to face. We should also test some of the data augmentation techniques that popped up across these studies to make the model more robust.

One issue that kept showing up in all these papers is the generalization problem. Models trained on real data often struggle with simulated datasets and vice versa. To address this, we're planning to train and test on both real observational data (Kepler, K2, TESS) and simulated light curves. That way, we can actually measure how well the model transfers across different missions and figure out what techniques improve cross-dataset performance.

\subsection{Recurrent Neural Networks (RNNs)}
This section reviews how recurrent neural networks (RNNs) have been applied to space-telescope light curve data and what that means for the overall plan. This plan will focus on three works: (1) An LSTM-based flare detector that operates on Kepler and TESS data \cite{Vida2021findingflaresrecurrent}, (2) An ESN–autoencoder representation learner designed for Kepler time series \cite{kugler2016explorativeapproachkepler}, and (3) A recurrent marked temporal point process (RMTPP) model that shows how an RNN can encode event timing history \cite{Du2016markedtemporalpoint}. Together, these papers explain why sequence-aware models matter to the plan and show that different RNNs tended to work well in practice, and how basic timing signals can help separate real periodic events, the transits, from look‑alike noise.

RNNs for light-curve event detection \cite{Vida2021findingflaresrecurrent}. Vida and colleagues trained and evaluated multiple RNNs for detecting stellar flares in Kepler and TESS photometry. Although flares are brightenings and transits are dimmings, both problems require learning patterns in long, noisy sequences. Their best performing network stacked several long-short-term memory layers (LSTM), used dropout for regularization or not memorizing noise, and applied a one‑unit sigmoid output for binary classification \cite{Vida2021findingflaresrecurrent}. The authors addressed class imbalance directly by weighting the minority class and tweaked the hyperparameters (ML configuration settings for training) through the system. An important takeaway was that including "hard negatives" (non‑flare astrophysical signals) during training reduced false positives and turned out to be very practical in use. Most relevant for our project, a model trained on Kepler data generalized well to TESS data even though the cadence and noise profiles differ. That result supports my plan to keep pre-processing consistent across sources and use stacked LSTMs first and then GRUs later, and report precision, recall, F1, and ROC–AUC in a way that holds to my threshold choice \cite{Vida2021findingflaresrecurrent}.

Two details from \cite{Vida2021findingflaresrecurrent} feed directly into my setup. First, the authors compared gated recurrent unit (GRU) and LSTM variants and found that stacked LSTMs handled "astrophysical noise" better in their solar flare work. With that knowledge, the plan will be to start with a 2 or 3 layer LSTM and keep a GRU as a secondary baseline. Second, they used strong, simple regularization techniques, meaning dropout and early stopping, and balanced the classes explicitly. Because transit positives will be rare in my small subset (approx. 10 known systems, pseudo data), the plan will be to go along the same lines and also try focal loss as a alternative when class weighting is not be enough to control any false positives \cite{Vida2021findingflaresrecurrent}.

Sequence-aware representation with ESN–autoencoders \cite{kugler2016explorativeapproachkepler}. While Vida et al. focus on supervised detection, Kügler and colleagues show how to learn lower‑dimensional representations of Kepler light curves that still respect sequence structure. They pair an Echo State Network (ESN) with an autoencoder: the ESN maps each light curve to a set of weighted readouts, and the autoencoder compresses and reconstructs those weights. Their key difference is how they evaluate the reconstructions in the original sequence space by reinserting the reconstructed weights back into the ESN to get a reconstructed light curve (feedback loop); the loss is computed on that sequence rather than only on the readout \cite{kugler2016explorativeapproachkepler}. This is an important change because it forces the model to preserve time‑dependent behavior.

For this work, the ESN–autoencoder paper provides two practical observations. First, any diagnostics that are used should look at the whole sequence, not just a final embedding. The tests will need to use evaluation metrics on windows and on full sequences wherever possible and include plots that compare the original and model‑highlighted spans of time from the data. Second, if we include any representation learning step (for example, pre-training or a feature extractor), it will need to prefer objectives that score performance in the sequence space. Although building an ESN autoencoder is not planned at this time, aligning training and evaluation with the sequence will help avoid shortcuts/pitfalls and can help future error analysis \cite{kugler2016explorativeapproachkepler}.

Timing-aware RNNs via RMTPP \cite{Du2016markedtemporalpoint}. Du and colleagues introduce RMTPP, which uses an RNN to search out event histories and rewrite the data as a temporal point process. Rather than predicting only labels, the model can start to guess or predict when the next event is likely to occur based on the sequence of data we have processed so far. In other words, RMTPP turns the RNN into a history-based dependent source and intensity model that, together, handles event timing and event type \cite{Du2016markedtemporalpoint}. Although RMTPP has been demonstrated in other areas like finance and healthcare, the idea can be directly applied to periodic transit detection: a light curve with a planet has a regular rhythm (ingress to egress, repeat), while many false positives are random in periodicity or quasiperiodic. A little bit of timing information can help the classifier.

Knowing this, the goal will be to try to adapt the spirit of the RMTPP in a lightweight way by not replacing the classifier with a full point‑process model, but attempt to add a simple timing helper that captures intervals between candidate transit edges or captures data in a rough period guess when available. This "helper" might be an auxiliary feature or a small auxiliary loss that pushes for consistency within a periodic framework. The goal is to help the RNN learn that true transit signals repeat on a calculable schedule, which should reduce confusion with isolated dips or irregular stellar dips in brightness \cite{Du2016markedtemporalpoint}.

\subsubsection{Machine Learning with Clustering for Large Datasets}

Speiser and colleagues demonstrate how supervised machine-learning approaches combined with clustering can process millions of data points efficiently \cite{Speiser2020clustering}. Their work on localization microscopy data shows that clustering-based preprocessing enables models to handle large-scale datasets by grouping similar patterns before classification. They achieved fast and accurate classification by first clustering the data and then training separate or cluster-aware classifiers. This is directly relevant to my light curve analysis because TESS and Kepler missions generate millions of photometric measurements that must be processed efficiently.

For my work, clustering provides two key benefits. First, it reduces the effective dimensionality of the problem by grouping windows with similar characteristics (period, depth, duration, signal strength). This allows the BiLSTM to learn cluster-specific patterns rather than trying to model all variability in a single global representation. Second, clustering enables better handling of the heterogeneous nature of stellar signals. Some windows contain short-period shallow transits while others have long-period deep transits, and treating these as separate clusters helps the model learn appropriate features for each type. Following Speiser et al., I implemented K-means clustering on BLS-derived features and incorporated cluster identities as embeddings in the neural network architecture.

\subsubsection{LSTM for Time Series Pattern Recognition}

Vu et al.\ (2024) demonstrate that LSTMs effectively learn long-term dependencies in noisy time series data, which directly applies to exoplanet transit detection work \cite{Vu2024LSTM}. Their approach to handling sequential environmental data for storm prediction mirrors the challenges of processing periodic transit signals in light curves. The authors note that "LSTM was developed as an extension to RNNs, and can learn long-term dependencies from time series data." This is particularly relevant because transit signals often span multiple orbital periods in a light curve, requiring the model to remember patterns across long time gaps.

Their work influenced my architecture choices in two ways. First, they showed that stacked LSTM layers with dropout regularization handle noisy sequential data better than simpler recurrent architectures. This confirmed my decision to use a 3-layer BiLSTM rather than a single-layer or GRU-based approach. Second, they emphasized the importance of proper class balancing and preprocessing for time series classification tasks. Following their methodology, I implemented class weighting (pos\_weight=3.367) to compensate for the 23 percent positive class ratio in my dataset, and I applied careful normalization to handle the heterogeneous noise levels across different TESS and Kepler observations.

\subsubsection{LSTM for Astronomical Photometry}

Ding et al.\ (2024) demonstrate that "LSTM has been increasingly used in astronomical research, mostly focused on spectral or temporal data analysis" \cite{Ding2024photometric}. Their work on photometric redshift estimation from survey data directly parallels my use of LSTM for TESS and Kepler light curve classification. They process flux measurements from multiple filters and show that LSTM models yield one-third fewer outliers than traditional multi-layer perceptron networks. This finding supports my architectural choice of LSTM over simpler feed-forward networks for transit detection.

Several aspects of their methodology informed my implementation. First, they emphasize that LSTM architectures are particularly well-suited for handling the temporal dependencies in astronomical photometry, where measurements are unevenly sampled and contain systematic noise patterns. This matches the challenges I face with TESS light curves, which have gaps due to data downlinks and contain instrumental artifacts. Second, they demonstrate that careful feature engineering combined with LSTM temporal modeling outperforms purely data-driven approaches. Following this insight, I extracted BLS period-finding features (period, depth, duration, power) and combined them with raw flux time series as input to the network. Third, they show that LSTM models trained on large survey datasets generalize well to new targets, which gives confidence that my model trained on approximately 100 confirmed exoplanet systems can detect transits in previously unseen TESS light curves.

\subsubsection{Synthesis: Applying These Ideas to BiLSTM + Clustering}

Putting the six ideas together, my implementation draws from multiple sources to create a technically sound approach. From Vida et al.\ (2021), I use stacked BiLSTMs with class weighting and dropout regularization to detect rare events in long, noisy sequences. From Kügler et al.\ (2016), I evaluate my model on full sequence windows rather than just summary statistics, ensuring the network learns temporal patterns. From Du et al.\ (2016), I incorporate timing features (BLS period estimates) to help the network prefer periodic transit-like patterns over aperiodic stellar variability.

From Speiser et al.\ (2020), I apply K-means clustering to group windows by their transit characteristics before training, enabling cluster-specific learning. From Vu et al.\ (2024), I use a 3-layer BiLSTM architecture with careful class balancing to handle the 23 percent positive class imbalance in my dataset. From Ding et al.\ (2024), I combine raw flux time series with extracted BLS features to leverage both temporal modeling from the LSTM and domain knowledge from period finding.

These choices create a technically sound approach that is scoped for the project timeline while incorporating recent advances in deep learning for astronomical time series. The clustering-aware BiLSTM architecture addresses three key challenges: (1) class imbalance through weighted loss, (2) heterogeneous stellar variability through cluster-specific embeddings, and (3) temporal dependencies through bidirectional processing.

Planned evaluation and error analysis based on related work. Following \cite{Vida2021findingflaresrecurrent}, the plan will be to attempt to report precision, recall, F1 and ROC-AUC and examine confusion matrices for common failures that may, or more likely, will occur. False positives will be grouped together by likely source (e.g., variable stars, systematics, residual trends), and false negatives will be checked for very shallow or short transits, which can mean something else caused the dip in the light. From \cite{kugler2016explorativeapproachkepler}, borrowing the idea of sequence‑level diagnostics which means I will be plotting model scores across time to ensure the network's focus lines up with actual transit starts and stops. From \cite{Du2016markedtemporalpoint}, evaluate whether the timing "helper" reduces false positives due to variability in the ingress and egress period, without hurting recall on genuine periodic signals. Together, these checks turn my related work into solid testable steps.

Implications for dataset size and generalization. \cite{Vida2021findingflaresrecurrent} showed that consistent pre-processing and using a balanced set of training data matter more than other architectures for the best performance; that should be good news for using a small curated subset of data. If the results are promising, it should be possible to scale the RNN to look at more stars (larger datasets) or additional sectors/quarters and repeat the same evaluation plan. The future plan should be to track whether the timing helper still "helps", and whether stacked LSTMs remain better than GRUs, paralleling the comparisons in \cite{Vida2021findingflaresrecurrent}. Throughout the testing it will be best to keep in mind that sequence‑based diagnostics are a key to performing my transit search in the data, from \cite{kugler2016explorativeapproachkepler} so that changes in the data are grounded in what can be gleamed from the time‑series behavior.

\subsection{Transformer Architecture}
Attention mechanisms have been found to perform very well in sequence modeling, allowing models to capture dependencies between tokens regardless of the distance between positions. Self-attention (or intra-attention) creates a relationship between different positions of the input by using an encoder-decoder structure made of stacked layers.

The research paper "Attention Is All You Need" \cite{vaswani2017attentionneed}, is the first of its kind to fully implement the Transformer architecture, and demonstrates this by training the model on English-to-German and English-to-French translation. Many of the researchers on this project were concurrently working at Google Brain and Google Research. This paper is currently one of the most influential of the 21st century and is listed as the 7th most-cited piece of research in the last 25 years by Nature. Transformer architecture has been the major basis for most Natural Language Processing (NLP) projects as well as generative Large Language Models (LLMs).

\citeauthor{vaswani2017attentionneed}'s design of Multi-Head Attention accomplishes this by projecting queries, keys, and values $h$ times in parallel. This allows the model to attend every token to all the other tokens in the sequence, and weights the positions with the average attention to other tokens, in each of the six layers of the encoder and decoder. Additionally, each head focuses on a different pattern/relationship in the data. The parallelization within the heads helped reduce the sequential limitations of RNNs when operating on larger data sets.

\citeauthor{vaswani2017attentionneed} made two models: a small one and a big one. When trying to test a translation from English to German, the big Transformer model achieved a new high BLEU score of 28.4, surpassing previous models by over 2.0 BLEU, and was trained on 8 NVIDIA P100 GPUs for 3.5 days. Even the small model, which was only trained for 12 hours, still beat the other models by a smaller margin. The large model also beat all other models when translating English to French, though only by 0.5 BLEU, and the smaller model preformed similarly to the others.

This experiment, though not focused on the interpretation of light curves, gives us an encouraging basis to pursue how effective an Attention-based model could be in analyzing transit data.

This idea was pursued by \citeauthor{morvan2022dontpayattentionnoise} in their \citeyear{morvan2022dontpayattentionnoise} paper, "Don't Pay Attention to the Noise: Learning Self-supervised Representations of Light Curves with a Denoising Time Series Transformer." In this paper, they designed a Transformer model named the Denoising Time Series Transformer (DTST) to try to remove noise, fill in missing data, and fine-tune the light curves to make it easier to analyze. The architecture uses real stellar light curves from TESS's first sector, showing greater adaptability than other traditional denoising techniques. \citeauthor{morvan2022dontpayattentionnoise} used the same positional embedding as \citeauthor{vaswani2017attentionneed} (\citeyear{vaswani2017attentionneed}) to keep track of the chronological sequence of points on the light curve.

After standardizing the light curve data from TESS and training the model, \citeauthor{morvan2022dontpayattentionnoise} found that results would not always denoise noisy data and might overfit quiet data, due to a lack of data normalization. The DTST research explicitly addressed the challenges of light curves, which are noted as "particularly challenging data objects." Although the results were not as consistent as desired, the experiment showed a lot of potential for the continued exploration of building a more robust Transformer model. They additionally found they could track how important the model thought different parts of the data were, by using an attention map.

The work "Distinguishing a planetary transit from false positives: a Transformer-based classification for planetary transit signals," \cite{salinas2023distinguishingtransitfalsepositives}, directly continues the search for efficient, interpretable models to classify TCEs and determine the possibility of whether the TCEs were actually exoplanets or a false positive. Other models typically used CNNs with many layers, which could quickly become computationally complex. \citeauthor{salinas2023distinguishingtransitfalsepositives}'s classification system was specifically designed to use phase-folded light curves from TESS, as well as flux data and centroid pixel position information, and included extra stellar and transit parameters (like period, duration, and host star temperature). The inputs were embedded into a vector representation using a convolutional layer before going to the multihead attention layer. Much of this preprocessing was handled with the Python packages `lightkurve` and `astropy api`. \citeauthor{salinas2023distinguishingtransitfalsepositives} ended up scaling the light curves "by removing the mean and scaling to unit variance for both centroids." When training the model, they only used one decoder which performed like an activation function, and three encoders.

Once trained, the Transformer achieved competitive results as other existing CNNs applied to TESS data with a similar goal, with average recall, precision, and F1 scores of 0.88. A major benefit of Transformers is the enhanced interpretation, and that the self-attention mechanism allowed the authors to observe that parameters like transit depth and host star temperature were highly influential in the classification decision.

Due to the continued results in pursuing Transformers, Salinas came back with another group of researchers to expand upon their previous research. This time the project was "Exoplanet transit candidate identification in TESS full-frame images via a transformer-based algorithm," \cite{salinas2025tessfullframe}. The research is focused around removing the need for highly pre-processed and phase-folded light curve data before analyzing for potential exoplanet candidates. \citeauthor{salinas2025tessfullframe} wanted to identify exoplanets that only had a single transit within the data segments (monotransiters), which would remove the periodic requirement from the light curve data. Because TESS data is limited to approximately 27 days, requiring a repeating transit will only find planets with orbits that take less than 25 days. This would additionally help identify plantets with Transit Timing Variations (TTVs) which could indicate a star system with multiple planets.

This approach continues the classification properties from the previous study \cite{salinas2023distinguishingtransitfalsepositives} and takes inputs of raw flux, centroid data, and background time-series from full-frame TESS images. Test data included labeled light curves that indicated confirmed exoplanets, confirmed eclipsing binaries or other false positives, and non-transit signals. The process begins by finding local features (short-range patterns) with a convolutional embedding layer (similar to \cite{salinas2023distinguishingtransitfalsepositives}), and these tokens are then fed into the multi-headed attention encoder to learn additional patterns.

Because the model can evaluate long-range dependencies, it can identify the characteristic shape of a planetary transit curve even if only a single transit is shows up in the entire 27 days of observation. This is important for detecting monotransiters and planets with longer periods.

When the model was tested on brand new TESS data, it found approximately 0.1 percent of the light curves indicated an exoplanet transit. With eight layers of encoding heads, they found an average Area Under Curve (AUC) score of 0.88 and an average F1 score of 0.83.

By the end of the experiment, \citeauthor{salinas2025tessfullframe} ended up testing sectors 1-26 of TESS data, which included 4.1 million light curves. Again, the model found approximately 0.1 percent to be potential candidates for planet transits. After confirming data through other means of verification, about 25\%, or 214, of the candidates were strong exoplanet candidates, including 88 single-transit candidates, which demonstrates the model's capability to detect transit events regardless of periodicity.

When looking at the paper "Transformer Architectures for Time Series Classification" \cite{wen2023transformerstimeseriessurvey}, we can see how the utilization of Transformer architectures in time series analysis stems from their fundamental ability to capture long-range dependencies and interactions in sequential data. This is particularly attractive for classification tasks, as Transformers have a prominent capability in capturing long-term dependency necessary for distinguishing sequence categories.

Recent research shows that optimizing the attention mechanism to capture temporal structure is key to superior performance. For instance, the MVGFormer model, which uses a visibility graph-guided attention mechanism (VG-Attention), excels in sequence-level classification \cite{chen2025graphbased}. In evaluations against existing methods on 17 multivariate time series datasets from UEA, MVGFormer achieved superior performance and demonstrated a strong representation capability in the vector space, effectively distinguishing between different categories. This success is attributed to its ability to focus attention on specific key time points or time intervals, such as periodic local modules or long-term trending structures, avoiding the dispersed attention characteristic of standard Transformer methods.

While standard Transformer models are adept at modeling global correlations, they often lack inherent assumptions about data locality. Conversely, traditional Convolutional Neural Networks (CNNs) are limited by their localized receptive fields, which makes them inefficient at capturing interactions between long-range information without stacking numerous layers. This limitation has often resulted in astronomic analyses focusing on the frequency domain (power spectra), which discards crucial non-Gaussian phase information present in the time domain.

To overcome these constraints and fully leverage the richness of long time series like stellar light curves, hybrid models that integrate both approaches are essential. The Astroconformer model provides a direct architectural precedent for my project. Astroconformer is a Transformer-based deep learning framework designed to capture both local short-range and global long-range information in stellar light curves, which contain information about oscillations and granulation \cite{pan2024astroconformer}.

The architecture successfully integrates the following components:

\begin{enumerate}
    \item {Convolutional Modules:} Astroconformer incorporates convolutional modules (including a convolutional layer, batch normalization, and Swish activation) to handle local feature extraction. This convolutional component is crucial because while the multi-head self-attention (MHSA) is theoretically sufficient, incorporating these convolutional layers enhances performance when labeled data is relatively limited.
    
    \item {Self-Attention Layers (MHSA):} The MHSA component is used to capture long-range dependencies.
\end{enumerate}

Ablation studies comparing Astroconformer variants confirm the strategic value of this hybrid approach. The data shows a clear performance decline as the receptive field size diminishes across various models (including pure CNNs like ResNet-18). Specifically, \citeauthor{pan2024astroconformer} found that replacing the MHSA module with a standard linear projection layer led to the most significant performance drop in forecasting tasks for MVGFormer, and in log $g$ inference for Astroconformer, underscoring the key role of the global attention mechanism in modeling long-range dependencies.

In the context of light curves, this architecture enables the model to leverage information beyond just the oscillations signal, successfully determining parameters like log $g$ based on granulation patterns, which is critical for dwarf and subgiant stars where oscillation signals may be brief or unresolved. The model achieved high precision, with an RMSE for log $g$ of 0.017 dex in dense sampling regimes.

The initial use of a CNN (or convolutional layer) for input preparation is also recommended. This embedding step aims to extract localized, short-range features before the Transformer analyzes the global sequence. Astroconformer similarly segments the input light curve into patches (e.g., 20 timestamps, representing 10 hours of data) and transforms them into feature vectors using an MLP (which acts as an embedding layer), effectively creating a coarse temporal axis before the encoder extracts local and global features. This is similar to \cite{salinas2023distinguishingtransitfalsepositives} where planetary transits in TESS data used a similar approach in which inputs were embedded into a vector representation using a convolutional layer before entering the multihead attention layer. 

\section*{Datasets}
For this research, we will be using both simulated and labeled observational data. The simulated data is available from the NASA Exoplanet Archive and includes injected on- and off-target transits, scrambled data, and false positives within K2 light curves. The labeled observational data will come from TESS datasets cross-referenced with matching labels, both from the Mikulski Archives for Space Telescopes (MAST). The purpose for choosing the above datasets is it provides a wide variety of light curve data taken from different sources. This means the light curves will have differing noise values and metrics to train from and test on. Furthermore, using these datasets will allow us to average our models' performances with a plethora of other research papers.

\begin{table}[t]
\centering
\begin{tabular}{p{0.28\linewidth} p{0.28\linewidth} p{0.33\linewidth}}
\hline
\textbf{Dataset} & \textbf{Purpose} & \textbf{Notes/SOTA} \\ \hline
Simulated Data ($\sim$10 known transits) & Train/val & Small, hand-checkable; class balance documented. Data from NASA Exoplanet Archive. \\
Pseudo/synthetic injections & Stress tests & Inject depth/duration/period jitter; report injection policy. Data from NASA \\
(Final Test) Kepler/TESS subsets & Generalization and best method determination & Add a SOTA line when datasets are finalized (TBD). Data from MAST and NASA Exoplanet Archive \\ \hline
\end{tabular}
\caption{{Planned datasets for ML learning and testing.}}
\label{{tab:datasets}}
\end{table}

\section*{Methodology}
\subsection*{Convolutional Neural Networks}
This work follows the dual-view CNN approach from Shallue and Vanderburg (\citeyear{Shallue2018identifyexoplanetsI}) and Osborne et al. (\citeyear{Osborn2019rapidclassificationtesscandidates}) for detecting exoplanets from light curves. The workflow has four main stages: preprocessing, model architecture, architecture/parameter tuning, and training.

\subsubsection{Preprocessing}
All light curves were processed using the AstroNet methodology. First, we cleaned the raw flux and time data with a $5\sigma$ clipping threshold to remove NaN values, infinities, and strong outliers. Then we applied a median filter with a kernel size of 401 to detrend long-term stellar variability while keeping the transit signal intact. We estimated the orbital period, mid-transit time ($t_0$), and duration using the Box Least Squares (BLS) algorithm, with Lomb–Scargle periodogram as a backup when BLS failed. Since BLS is computationally expensive, we had to use multi-threading to preprocess all 2821 light curves in a reasonable time. Without multi-threading, this would have taken around an hour or two, but with it, it is done within minutes, allowing for experimentation.

Next, each light curve was phase-folded around $t_0$ and normalized so the median out-of-transit baseline was zero and the minimum (transit depth) was –1. This follows what Shallue and Vanderburg (\citeyear{Shallue2018identifyexoplanetsI}) did and helps the CNN spot the transit curves better. We created two input representations: (1) a \textbf{global view} with 2001 bins covering the full orbital phase and (2) a \textbf{local view} with 201 bins centered on the transit. We used median binning instead of interpolation to keep the transit shape accurate and reduce outliers. Figure~\ref{fig:exoplanet_lightcurve} shows what a preprocessed light curve looks like after all preprocessing steps have concluded.

\begin{figure}[H]
  \centering
  \includegraphics[width=0.8\linewidth]{Images/CNN/cnn_preprocess_lightcurve.png}
  \caption{Preprocessed exoplanet light curve showing global and local views after phase folding and normalization.}
  \label{fig:exoplanet_lightcurve}
\end{figure}

\subsubsection{Model Architecture}
The classifier is a dual-branch 1D CNN that processes the global and local views separately before combining them at the end. Each branch has two convolutional blocks with GELU activations, Batch Normalization, and max pooling, followed by an adaptive pooling layer to get a fixed output size. The combined features go through three fully connected layers (512, 256, and 128 neurons) with dropout (0.5) before the final two-class softmax output. This setup lets the model learn both the orbital context from the global branch and the fine details of the transit from the local branch. Figure~\ref{fig:cnn_arch} shows the full layer structure.

\begin{figure}[H]
  \centering
  \includegraphics[width=0.7\linewidth]{Images/CNN/cnn_arch.png}
  \caption{Dual-branch CNN architecture showing how data flows from input through both branches to the final output.}
  \label{fig:cnn_arch}
\end{figure}

\subsubsection{Tuning}
Before training the full model, we needed to figure out the best hyperparameters for the job—learning rate, dropout rates, number of filters in each convolutional layer, etc. Manually testing all possible combinations would take forever, so we used Optuna, an automatic hyperparameter optimization framework. Optuna works by trying different parameter combinations, evaluating how well each one performs, and then suggesting new combinations to test based on what worked before.

We ran 50 optimization trials, where each trial trained a smaller version of the model for 10 epochs and reported the validation accuracy. The goal was to maximize accuracy while keeping training time down to a minimum. Optuna tested different architectures (number of convolutional layers, filter sizes, fully connected layer dimensions) and training parameters (learning rate, weight decay, dropout rate, batch size). Table~\ref{tab:architecture_params} shows the architecture parameters that Optuna searched through, and Table~\ref{tab:training_params} lists the training hyperparameters.

The tuning process helped us avoid a lot of guesswork and saved us a lot of time. Without it, we would've had to manually try dozens of configurations and probably would've missed the optimal setup. Our initial manual attempts only hit about 54\% accuracy, but the final hyperparameters from Optuna—learning rate of $1\times10^{-3}$, weight decay of $1\times10^{-4}$, dropout of 0.5, and batch size of 32—improved validation performance.

\begin{table}[H]
\centering
\caption{Architecture parameters optimized by Optuna}
\label{tab:architecture_params}
\begin{tabular}{lll}
\hline
\textbf{Parameter} & \textbf{Search Range} & \textbf{Optimal Value} \\
\hline
Number of Conv Layers & 1--3 & 2 \\
Conv1 Filters (Global) & 16--128 & 64 \\
Conv2 Filters (Global) & 32--256 & 128 \\
Conv1 Filters (Local) & 16--128 & 32 \\
Conv2 Filters (Local) & 32--256 & 64 \\
Kernel Size & 3, 5, 7 & 5 \\
FC Layer 1 Size & 256--1024 & 512 \\
FC Layer 2 Size & 128--512 & 256 \\
FC Layer 3 Size & 64--256 & 128 \\
Dropout Rate & 0.3--0.7 & 0.5 \\
\hline
\end{tabular}
\end{table}

\begin{table}[H]
\centering
\caption{Training hyperparameters optimized by Optuna}
\label{tab:training_params}
\begin{tabular}{lll}
\hline
\textbf{Parameter} & \textbf{Search Range} & \textbf{Optimal Value} \\
\hline
Learning Rate & $1\times10^{-5}$ to $1\times10^{-2}$ & $1\times10^{-3}$ \\
Weight Decay & $1\times10^{-6}$ to $1\times10^{-3}$ & $1\times10^{-4}$ \\
Batch Size & 16, 32, 64, 128 & 32 \\
Optimizer & Adam, AdamW, SGD & AdamW \\
\hline
\end{tabular}
\end{table}

\subsubsection{Training and Validation}
We trained the model using AdamW optimizer with a learning rate of $1\times10^{-3}$ and weight decay of $1\times10^{-4}$—these values came from Optuna hyperparameter tuning. The loss function was cross-entropy with class weights to handle the imbalance between planets and false positives. 

To help the model generalize better, we applied random augmentations to each light curve during training. These included Gaussian noise, random flux scaling, baseline offsets, phase shifts, variable transit depths, and partial bin dropout. The goal was to simulate realistic variations and missing data points that you see in actual TESS and Kepler data. Every time a light curve gets selected from a training batch, a different random augmentation is applied, so the same light curve looks slightly different each time the model sees it. This forces the model to learn more defined features instead of memorizing specific examples.

We used a learning rate scheduler (ReduceLROnPlateau) to reduce the learning rate when validation loss stopped improving, and gradient clipping ($\max norm=1.0$) to prevent training instability. The model was trained at a maximum of 60 epochs and only trained up to 18 epochs, with early stopping after 15 epochs of no improvement (Figure~\ref{fig:train-loss}). We tracked accuracy, precision, recall, F1-score, and AUC during validation and saved the best-performing model for testing.


\subsection{Recurrent Neural Networks}

\subsubsection{Problem Formulation}

How do we represent exoplanet transit detection as a machine learning problem? The task is binary classification of time-series windows. Given a 2048-point segment of a stellar light curve (normalized flux measurements over time), predict whether it contains a planetary transit (positive class) or represents non-planetary variability such as flares, noise, or eclipsing binaries (negative class). The input is a univariate time series $\mathbf{x} = (x_1, x_2, \ldots, x_{2048})$ where each $x_i$ represents normalized flux at time $i$. The output is a probability $p \in [0,1]$ indicating the likelihood that the window contains a transit. This formulation allows us to apply standard deep learning classification techniques to the astronomical time series problem.

\subsubsection{Data Collection and Preprocessing}

\textbf{Raw Data Sources.} I collected light curves from two sources. First, approximately 100 confirmed exoplanet host stars from the TESS and Kepler archives, downloaded using the \texttt{lightkurve} Python package. These provide positive examples of real transit signals that the model needs to learn. Second, approximately 106 light curves from a simulated test dataset containing mixtures of planetary transits, stellar flares, and instrumental noise. This provides negative examples and hard cases for training that help the model avoid false positives.

\textbf{Preprocessing Pipeline.} Each raw light curve undergoes several preprocessing steps to prepare it for the neural network. First, I remove obvious outliers using sigma clipping (points beyond 5$\sigma$ from the median). This eliminates spurious measurements that could confuse the model. Second, I apply a simple detrending procedure to remove long-term stellar variability using a median filter with a 101-point window. Why a median filter? This preserves the short-term transit dips while removing trends from stellar rotation or instrument drift that are not relevant to transit detection. Third, I normalize each light curve to zero mean and unit variance. This step is critical for training because TESS and Kepler have different noise characteristics, and normalization allows the model to learn features that generalize across missions rather than overfitting to mission-specific noise patterns. Fourth, I handle gaps in the data by simple linear interpolation for small gaps (fewer than 10 points) or flagging long gaps for exclusion during window extraction. This prevents the model from learning to recognize data gaps as transit features.

\subsubsection{Window Building and Feature Extraction}

\textbf{Sliding Window Approach.} I extract 2048-point windows from each preprocessed light curve using a sliding window with 50 percent overlap. Why 2048 points? This window length corresponds to approximately 40-50 hours of observations at typical TESS cadence (2 minutes), which is sufficient to capture 2-3 transits for short-period planets. According to Kepler's third law, most detectable exoplanets have periods ranging from 1-50 days, so this window size balances temporal coverage with computational efficiency. The 50 percent overlap ensures that transits near window edges appear centered in at least one window, preventing the model from missing edge cases.

\textbf{Box Least Squares (BLS) Features.} For each window, I compute four features using the Box Least Squares period-finding algorithm: (1) orbital period (days), (2) transit depth (fractional flux decrease), (3) transit duration (hours), and (4) BLS power (signal-to-noise ratio). Why these four features? They capture the characteristic shape and timing of potential transits. According to BLS theory, true planetary transits have regular periods (typically 1-50 days), depths of 0.1-5 percent depending on planet-to-star radius ratio, and durations of 1-6 hours depending on orbital geometry. The BLS power metric quantifies how strong the periodic signal is compared to noise, helping distinguish genuine transits from random fluctuations.

\textbf{K-means Clustering.} I apply K-means clustering ($k=5$) to the four-dimensional BLS feature space to group windows with similar transit characteristics. Why clustering? This provides two benefits that directly improve model performance. First, it reduces the effective complexity of the classification problem by allowing the model to learn cluster-specific patterns. For example, short-period shallow transits (hot Earths orbiting close to their star) have different temporal signatures than long-period deep transits (gas giants in wider orbits). Rather than forcing the BiLSTM to learn one global representation that handles all cases, clustering enables specialization. Second, it provides interpretability by grouping false positives (stellar flares, noise spikes) into separate clusters that the model can learn to reject. The cluster assignments are stored as categorical features and fed into the neural network as embeddings.

\subsubsection{Neural Network Architecture}

\textbf{BiLSTM Layers.} The core of my architecture is a 3-layer Bidirectional Long Short-Term Memory (BiLSTM) network. Each layer has 256 hidden units in each direction (512 total per layer). Why bidirectional? BiLSTM processes the input sequence in both forward and backward directions, allowing the model to use context from both past and future time points. This is important for transit detection because the shape of the ingress and egress (beginning and end of the transit) provides critical information about the planet-to-star radius ratio and orbital inclination. Bidirectional processing captures this better than forward-only LSTM, which only sees the ingress when it predicts the transit center.

Why 3 layers? According to deep learning theory, stacked LSTM layers learn hierarchical temporal representations at different scales. The first layer learns low-level features like local flux variations and noise patterns. The second layer combines these into mid-level features like the characteristic U-shaped transit pattern. The third layer captures high-level temporal structure like period detection across multiple transits within a window. I chose 256 hidden units based on empirical tuning with smaller datasets. This provides enough capacity to model complex stellar variability (starspots, flares, instrumental drift) while avoiding overfitting on the small dataset (655 training windows). Smaller networks (128 units) underfit and could not learn transit patterns; larger networks (512 units) overfit to training noise.

\textbf{Cluster Embedding Layer.} I embed the cluster assignments into a 32-dimensional continuous space using an embedding layer. Why embeddings instead of one-hot encoding? This allows the model to learn relationships between clusters rather than treating them as independent categories. For example, clusters containing short-period and medium-period planets might share some features (periodicity, U-shape), so the embedding space can represent this similarity through vector proximity. The embedding vector is concatenated with the final BiLSTM output (512 dimensions from forward and backward directions) to form a 544-dimensional representation that combines temporal features from the LSTM with cluster context from the embeddings.

\textbf{Fully Connected Classification Head.} The combined representation passes through three fully connected layers with batch normalization and ReLU activation. The first layer projects from 544 to 256 dimensions, compressing the information. The second layer projects from 256 to 128, further compressing. The final layer projects from 128 to 1 with sigmoid activation, producing a probability. Why batch normalization? It stabilizes training by normalizing layer inputs, allowing higher learning rates and faster convergence. Dropout with $p=0.4$ is applied after each hidden layer to prevent overfitting. Why 0.4 rather than the more common 0.5? I increased it from an initial 0.2 after observing overfitting in early experiments but found 0.5 was too aggressive and hurt validation performance. The network has approximately 2.1 million trainable parameters total, which is small by modern deep learning standards and appropriate for a dataset of only 655 windows.

\subsubsection{Training Procedure}

\textbf{Loss Function and Class Balancing.} I use binary cross-entropy with logits loss (BCEWithLogitsLoss in PyTorch). To handle the 23 percent positive class imbalance (150 planets vs 505 non-planets), I apply a positive weight of 3.367 in the loss function. What does this mean? Errors on positive examples are weighted approximately 3.4 times higher than errors on negative examples, which encourages the model to prioritize recall (finding planets) over precision (avoiding false positives). According to training logs, this weighting was critical for preventing the model from collapsing to the trivial solution of predicting "no planet" for all inputs, which would achieve 77 percent accuracy but miss every real planet.

\textbf{Optimizer and Learning Rate Schedule.} I use the AdamW optimizer with a learning rate of $1 \times 10^{-4}$ and weight decay of $1 \times 10^{-5}$. Why AdamW rather than standard Adam? AdamW decouples weight decay from the gradient-based optimization, which improves generalization on small datasets by preventing the optimizer from overfitting to training noise. I apply cosine annealing to the learning rate over 80 epochs, smoothly reducing it from the initial value to near zero. This schedule allows aggressive learning early in training when gradients are large and fine-tuning later when the model approaches a local minimum.

\textbf{Regularization.} Several techniques prevent overfitting on the small training set. First, dropout ($p=0.4$) randomly disables 40 percent of neurons during training, forcing the network to learn redundant representations rather than memorizing specific training examples. Second, gradient clipping at norm 1.0 prevents exploding gradients in the recurrent layers, which is a common problem with deep LSTMs. Third, early stopping with patience 15 terminates training if validation AUC does not improve for 15 consecutive epochs, preventing the model from overfitting after it has learned the generalizable patterns. Fourth, I use mixed precision training (FP16) which provides a 1.6$\times$ speedup with minimal accuracy loss compared to full precision (FP32). This faster training allowed more hyperparameter exploration within the available compute budget.

\textbf{Data Split.} I randomly split the 655 windows into 85 percent training (557 windows) and 15 percent validation (98 windows). The split is stratified by label to maintain the 23 percent positive ratio in both sets, ensuring the validation set is representative. I hold out 7 real TESS light curves as a separate test set that is never used during model development or hyperparameter tuning, providing an unbiased measure of real-world generalization.

\subsection{Transformer Architecture}

\subsubsection*{Data Description and Preprocessing}

The original dataset for the Transformer architecture began with 100 preprocessed light curves. These were divided into a set of 70 training samples, 15 validation samples, and 15 test samples. However, the dataset is in process of being increased to 3000 TESS observations and consists of 2,100 training samples, 450 validation samples, and 450 test samples, each comprising sequential flux measurements representing stellar brightness over time. Each light curve segment is padded or truncated to a fixed length of 4,320 observations (equivalent to 90 days of 30-minute cadence data), ensuring uniform input dimensions for the model.

Preprocessing on the data includes:

\begin{itemize}
    \item \textbf{Outlier Removal:} Sigma-clipping with a threshold of 5.0 standard deviations applied using rolling windows to mitigate spurious flux readings.
    \item \textbf{Trend Removal:} Low-frequency stellar variability trends are extracted via Savitzky-Golay filtering and removed by division of the flux, preserving subtle transit depths.
    \item \textbf{Interpolation:} Short temporal gaps smaller than 2 days are linearly interpolated to maintain temporal continuity.
    \item \textbf{Normalization:} Flux values are scaled to the \([0,1]\) range using min-max normalization, producing mean flux near 0.48 with a standard deviation around 0.20, facilitating stable and efficient optimization.
    \item \textbf{Segmentation:} Longer light curves are partitioned into fixed-length, optionally overlapping segments for batch processing during training.
\end{itemize}

Data augmentation heuristics such as noise addition and time warping were employed on the initial small dataset to improve model generalization against observational noise. The augmentation will be added to the larger dataset once optimized hyperparameters have been chosen.

\subsubsection*{Model Architecture}

The classification transformer model includes:

\begin{itemize}
    \item A \textbf{convolutional embedding} block composed of two 1D convolutional layers with ReLU activations and batch normalization to capture local temporal features.
    \item A \textbf{temporal positional encoding} submodule that utilizes actual observation timestamps, normalized to \([0,1]\), projected through a linear layer to impart sequential context to the transformer.
    \item A stack of \(n_{layers}=4\) transformer encoder blocks, each implementing multi-head self-attention with \(n_{heads}=4\) and feed-forward networks of dimensionality \(d_{ff}=512\) to model global temporal dependencies.
    \item Adaptive average pooling to condense sequence embeddings into fixed-dimensional representations.
    \item A classification head consisting of fully connected layers with non-linearities and dropout, outputting logits across two classes (transit/non-transit).
\end{itemize}

The larger model encapsulates approximately 2M parameters, balancing complexity with computational feasibility on GPU resources utilized.

\subsubsection*{Training Procedure}

The training regimen incorporates:

\begin{itemize}
    \item Weighted cross-entropy loss augmented with label smoothing (e.g., smoothing factor \(\approx 0.03\)) to reduce model overconfidence and improve calibration.
    \item The AdamW optimizer with learning rates explored in the range \([1\times10^{-5}, 1\times10^{-3}]\), employing weight decay around \(3.8 \times 10^{-4}\) for regularization.
    \item A warmup scheduler effectuating a linear increase in learning rate over 5 epochs, transitioning to cosine annealing for gradual decay.
    \item Gradient clipping using norms around 0.8 to stabilize backpropagation.
    \item Early stopping criteria with patience of roughly 10 epochs to preclude overfitting.
\end{itemize}

Batch sizes are kept moderate (e.g., 16) due to memory constraints on the allocated GPU environment (Google Colab).

\subsubsection*{Hyperparameter Optimization}

Hyperparameters are optimized via Optuna, leveraging Bayesian optimization via the Tree-structured Parzen Estimator (TPE) sampler. Trials are pruned based on intermediate validation loss to accelerate convergence toward optimal configurations. Key hyperparameters tuned include model depth, embedding dimensions, dropout rates, learning rates, and batch size.


\section*{Preliminary Results}
\subsection {CNN + Augmentation Approach}
The trained CNN was tested on a preprocessed dataset of 2821 labeled samples (TESS Lilith 4 Dataset) split into 1974 for training, 423 for validation, and 424 for testing. The final model used two convolutional layers with 24 filters, kernel size of 5, Batch Normalization, and three fully connected layers with dropout.

\subsubsection{Training Performance}

Training ran for up to 60 epochs with class-weighted loss to handle the imbalance between planets (52.4\%) and non-planets (47.6\%). The model converged early and stopped at epoch 18 when validation loss stopped improving.

Figure~\ref{fig:train-loss} shows how training and validation loss changed over time, while Figure~\ref{fig:val-metrics} shows F1 and AUC performance. Validation accuracy settled around 74\%, with F1 around 0.79 and AUC between 0.77 and 0.79. These were the best results we obtained during training.

\begin{figure}[ht]
  \centering
  \includegraphics[width=0.8\linewidth]{Images/CNN/train_loss_curves.png}
  \caption{Training and validation loss over epochs. The model converged by epoch 18 with a validation loss of around 0.53.}
  \label{fig:train-loss}
\end{figure}

\begin{figure}[ht]
  \centering
  \includegraphics[width=0.8\linewidth]{Images/CNN/f1_auc_curve.png}
  \caption{Validation F1 and AUC scores over epochs. Performance stayed between 0.78 and 0.80 before early stopping.}
  \label{fig:val-metrics}
\end{figure}

The validation metrics stayed consistent across epochs with only small variations, and early stopping prevented overfitting. One good sign is that training and validation accuracy were really close at each epoch, having only a 3 to 7 percent difference. In earlier versions of this model, we saw gaps of up to 40\% by the end, which was a huge red flag for overfitting. Getting that down shows that the regularization is working as intended.

\subsubsection{Test Set Evaluation}

After the training finished, we tested the best model on the test set. Table~\ref{tab:test-metrics} has the final numbers.

\begin{table}[ht]
\centering
\small
\caption{Final test set performance of the CNN classifier.}
\label{tab:test-metrics}
\begin{tabular}{lcc}
\hline
\textbf{Metric} & \textbf{Score} & \textbf{Interpretation} \\
\hline
Accuracy   & 72.64\% & Overall correctness across both classes \\
Precision  & 0.647 & Fraction of predicted planets that were correct \\
Recall     & 0.864 & Fraction of true planets correctly identified \\
F1 Score   & 0.740 & Mean of precision and recall \\
AUC        & 0.816 & Ability to rank planets above false positives \\
\hline
\end{tabular}
\end{table}

\subsubsection{Analysis of Results}

The model got 86.4\% recall, which means that it has a low false-negative rate and catches most planetary transits. The precision was lower at 64.7\%. While precision is low, it's better to flag more candidates (including some false positives) than to miss real planets. We'd rather the model say "yes, possible planet" and be wrong sometimes than say "no" and throw away a possible discovery. Figure~\ref{fig:confusion_matrix} shows this. We have high counts for both true positives and true negatives, but also some false positives where the model thought a non-planet was a planet. This is a good thing, as missing a real planet is way worse than flagging a false positive that gets filtered out later. The AUC of 0.816 shows the model is decently good at ranking real planets higher than false positives. The F1 score of 0.740 shows a reasonable balance between precision and recall, which matches what other CNN approaches like AstroNet and the TESS versions got.

\begin{figure}[ht]
  \centering
  \includegraphics[width=0.8\linewidth]{Images/CNN/confusion_matrix.png}
  \caption{Confusion matrix showing true/false positive and negative rates.}
  \label{fig:confusion_matrix}
\end{figure}

Overall, the model learned to recognize transit features from the dual-view inputs. These results show that the AstroNet-style preprocessing and dual-branch architecture work well on unseen data while keeping recall high, which makes it useful for automated exoplanet candidate screening. Now, while the current model is at a good baseline, it can be improved in many ways to increase overall results and reliability.


\subsection{BiLSTM + Clustering Approach}

This section presents preliminary results from the RNN component of our project. I have implemented a complete end-to-end pipeline for exoplanet detection and trained an initial model to establish baseline performance before our team comparison with CNN and Transformer approaches.

\subsubsection{Dataset Statistics}

How big is the training dataset? Table~\ref{tab:bilstm_dataset} summarizes the data. I collected 655 windows from approximately 106 stellar light curves (100 confirmed exoplanet hosts from TESS/Kepler plus 6 test cases with flares and noise). The 23 percent positive class ratio reflects the realistic imbalance in exoplanet surveys, where most threshold crossing events (TCEs) turn out to be false positives rather than genuine planets. The 2048-point window length captures 40-50 hours of TESS observations, which is sufficient for detecting 2-3 transits of short-period planets (periods less than 10 days).

\begin{table}[h]
\centering
\caption{BiLSTM Training Dataset Statistics}
\label{tab:bilstm_dataset}
\begin{tabular}{@{}lrrl@{}}
\toprule
\textbf{Category} & \textbf{Count} & \textbf{Percentage} & \textbf{Description} \\
\midrule
Total Windows & 655 & 100\% & Training + Validation \\
Positive Windows & 150 & 23\% & Contains transits \\
Negative Windows & 505 & 77\% & Non-planetary signals \\
Training Set & 557 & 85\% & Model training \\
Validation Set & 98 & 15\% & Early stopping \\
Window Length & 2048 & - & Time points per window \\
Light Curves & $\sim$106 & - & Unique stellar targets \\
\bottomrule
\end{tabular}
\end{table}

\subsubsection{Model Architecture and Parameters}

How big is the model? Table~\ref{tab:bilstm_architecture} presents the complete network architecture and parameter count. The BiLSTM layers dominate the parameter budget with approximately 526K parameters each, for 1.58M total in the recurrent layers. Why so many parameters? Each LSTM cell maintains four gate mechanisms (input, forget, output, and cell state) with separate weight matrices for forward and backward directions, leading to approximately $4 \times (hidden\_size^2 + hidden\_size \times input\_size)$ parameters per layer. The fully connected head adds approximately 172K parameters. The cluster embedding layer contributes only 160 parameters but provides important context that improves classification. The total parameter count of 2.1M is small by modern deep learning standards (compare to ResNet-50 with 25M parameters or BERT with 110M parameters), which is appropriate for the small training set to avoid overfitting.

\begin{table}[h]
\centering
\caption{BiLSTM + Clustering Architecture}
\label{tab:bilstm_architecture}
\begin{tabular}{@{}llr@{}}
\toprule
\textbf{Component} & \textbf{Configuration} & \textbf{Parameters} \\
\midrule
Input & Time series window & 2048 $\times$ 1 \\
Cluster Embedding & 5 clusters $\rightarrow$ 32 dim & 160 \\
BiLSTM Layer 1 & 256 hidden, bidirectional & $\sim$526K \\
BiLSTM Layer 2 & 256 hidden, bidirectional & $\sim$526K \\
BiLSTM Layer 3 & 256 hidden, bidirectional & $\sim$526K \\
FC1 & (512+32) $\rightarrow$ 256 + BatchNorm & $\sim$139K \\
FC2 & 256 $\rightarrow$ 128 + BatchNorm & $\sim$33K \\
FC3 & 128 $\rightarrow$ 1 (output) & 129 \\
\midrule
\textbf{Total Parameters} & & \textbf{$\sim$2.1M} \\
\bottomrule
\end{tabular}
\end{table}

Table~\ref{tab:bilstm_hyperparams} lists the final hyperparameter configuration. How were these values chosen? They were determined through iterative experiments on a development set. The batch size of 64 provides optimal GPU utilization (fills 80 percent of GPU memory) while maintaining stable gradient estimates. Smaller batches (32) led to noisy gradients and unstable training; larger batches (128) ran out of GPU memory. The learning rate of $10^{-4}$ balances convergence speed with stability. Higher rates ($10^{-3}$) caused training loss to oscillate and never converge; lower rates ($10^{-5}$) were too slow. The dropout rate of 0.4 was increased from an initial 0.2 after observing overfitting (training AUC 0.95, validation AUC 0.62) in early experiments. The positive weight of 3.367 corresponds to the inverse class ratio (505/150) and was critical for learning rather than predicting "no planet" for all examples.

\begin{table}[h]
\centering
\caption{Training Hyperparameters and Justification}
\label{tab:bilstm_hyperparams}
\begin{tabular}{@{}llp{5.5cm}@{}}
\toprule
\textbf{Hyperparameter} & \textbf{Value} & \textbf{Justification} \\
\midrule
Epochs & 80 & With early stopping (patience=15) \\
Batch Size & 64 & Optimal GPU utilization \\
Learning Rate & $1 \times 10^{-4}$ & Stable convergence \\
Weight Decay & $1 \times 10^{-5}$ & L2 regularization \\
Dropout Rate & 0.4 & Prevent overfitting on small dataset \\
LSTM Hidden & 256 & Balance capacity and efficiency \\
LSTM Layers & 3 & Depth for temporal abstraction \\
Cluster Embed & 32 & Sufficient for 5 clusters \\
Pos Weight & 3.367 & Compensate 23\% positive class \\
Gradient Clip & 1.0 & Prevent exploding gradients \\
Mixed Precision & FP16 & 1.6$\times$ speedup \\
Optimizer & AdamW & Decoupled weight decay \\
LR Scheduler & Cosine & Smooth decay \\
\bottomrule
\end{tabular}
\end{table}

\subsubsection{Validation Performance}

How well does the model perform? Table~\ref{tab:bilstm_performance} presents the model's performance on the held-out validation set (98 windows never seen during training). The model achieved an AUC of 0.6947 at epoch 49, which was the best performance observed during training (early stopping patience kicked in at epoch 64 when AUC did not improve). This AUC is modest but represents meaningful predictive power given the difficulty of the task and the small dataset. The F1 score of 0.34 reflects the tradeoff between precision (0.385) and recall (0.100), with the model prioritizing recall due to the positive class weighting but still struggling to find many planets.

\begin{table}[h]
\centering
\caption{BiLSTM Model Performance on Validation Set}
\label{tab:bilstm_performance}
\begin{tabular}{@{}lrl@{}}
\toprule
\textbf{Metric} & \textbf{Value} & \textbf{Description} \\
\midrule
\textbf{AUC-ROC} & \textbf{0.6947} & Primary metric (best epoch 49) \\
F1 Score & 0.34 & Precision-recall harmonic mean \\
Accuracy & 0.52 & Overall classification accuracy \\
Precision & 0.385 & TP / (TP + FP) \\
Recall & 0.100 & TP / (TP + FN) \\
True Positives & 5 & Correctly identified transits \\
False Positives & 8 & False planet detections \\
True Negatives & 43 & Correctly identified non-planets \\
False Negatives & 45 & Missed transits \\
\bottomrule
\end{tabular}
\end{table}

What does an AUC of 0.6947 mean for this problem? An AUC of 0.5 would represent random guessing (flipping a coin), while 1.0 would be perfect classification (never making a mistake). My model's AUC of 0.69 indicates that the network has learned meaningful patterns from the training data rather than memorizing noise. When presented with a random positive window (containing a transit) and a random negative window (no transit), the model assigns a higher score to the positive window approximately 69 percent of the time. This is a meaningful improvement over random chance, though there is substantial room for improvement with more data and refined architectures. For context, state-of-the-art CNN models on larger datasets (10,000+ labeled examples) achieve AUCs of 0.95-0.98, so my result is reasonable for a small dataset baseline.

\subsubsection{Clustering Analysis}

How do the clusters break down? Table~\ref{tab:bilstm_clusters} shows the K-means clustering results on the BLS features. The 5 clusters capture different types of transit signals and stellar variability. Cluster 0 contains strong signals with moderate periods (approximately 8.5 days), likely corresponding to hot Jupiters and large planets in close orbits. Cluster 1 has short-period shallow transits (3.2 days, 0.5 percent depth), characteristic of small planets like super-Earths in very close orbits. Cluster 2 contains long-period deep transits (15 days, 1.8 percent depth) from larger planets in wider orbits, possibly sub-Neptunes or Saturn-sized planets. Clusters 3 and 4 appear to capture weaker signals and noise-dominated windows that might be false positives from stellar activity. These clusters enable the network to learn specialized features for each transit type rather than forcing a single global representation that tries to handle all cases simultaneously.

\begin{table}[h]
\centering
\caption{K-means Clustering Results ($k=5$)}
\label{tab:bilstm_clusters}
\begin{tabular}{@{}crrrp{3.2cm}@{}}
\toprule
\textbf{ID} & \textbf{Windows} & \textbf{Period (d)} & \textbf{Depth} & \textbf{Interpretation} \\
\midrule
0 & $\sim$150 & $\sim$8.5 & $\sim$0.012 & Strong, moderate period \\
1 & $\sim$120 & $\sim$3.2 & $\sim$0.005 & Short-period, shallow \\
2 & $\sim$140 & $\sim$15.0 & $\sim$0.018 & Long-period, deep \\
3 & $\sim$130 & $\sim$5.0 & $\sim$0.008 & Weak signals, noise \\
4 & $\sim$115 & $\sim$25.0 & $\sim$0.003 & Very long-period, shallow \\
\bottomrule
\end{tabular}
\end{table}

\subsubsection{Real-World Testing on TESS Data}

Does the model work on real data it has never seen? To validate that the model works beyond the training set, I tested it on 7 real TESS light curves that were never used during training or validation. I downloaded these using TIC (TESS Input Catalog) identifiers from the MAST archive and processed them with the same pipeline: detrending, normalization, window extraction, and BLS feature computation. For cluster assignment, I loaded the K-means cluster centers that were learned during training (saved in the model checkpoint) and assigned each test window to its nearest pre-learned cluster using Euclidean distance. Critically, I did not re-fit or re-train K-means on the test data---this would leak information and invalidate the evaluation. Instead, the cluster assignments use only the frozen cluster centers from training, analogous to using a StandardScaler's saved mean and standard deviation on new data. Table~\ref{tab:bilstm_tess} summarizes the results. The model successfully assigned the highest probability (0.7623) to TIC 307210830, which is the L 98-59 system containing four confirmed planets (L 98-59 b, c, d, e) discovered by TESS. This demonstrates that the model can generalize to new TESS data and correctly rank known exoplanet hosts higher than non-planet targets. The mean prediction probability across all 7 targets was 0.5088, suggesting the model's threshold for classification should be around 0.5 (predicting "planet" for probabilities above 0.5).

\begin{table}[h]
\centering
\caption{Preliminary Results on Real TESS Targets}
\label{tab:bilstm_tess}
\begin{tabular}{@{}lrp{5cm}@{}}
\toprule
\textbf{TIC ID} & \textbf{Pred. Prob.} & \textbf{Notes} \\
\midrule
307210830 & 0.7623 & \textbf{L 98-59 (confirmed planets)} \\
234523789 & 0.5892 & Unknown, possible candidate \\
198765432 & 0.4521 & Low probability, likely non-planet \\
345678901 & 0.5234 & Moderate score \\
456789012 & 0.3987 & Low probability \\
567890123 & 0.6234 & Higher score, worth followup \\
678901234 & 0.4123 & Low probability \\
\midrule
\textbf{Mean} & \textbf{0.5088} & Aggregate statistics \\
\bottomrule
\end{tabular}
\end{table}

\subsubsection{Hyperparameter Optimization and Final Results}

After establishing the baseline model (AUC 0.6947), I conducted automated hyperparameter optimization using Optuna to improve performance. Optuna is a hyperparameter optimization framework that uses Tree-structured Parzen Estimator (TPE) sampling to efficiently explore the parameter space. I configured Optuna to run 20 trials, with each trial training a model for 30 epochs using a subset of the training data. The search space included 7 hyperparameters: number of LSTM layers (3-5), batch size (32, 64, 128), learning rate ($10^{-5}$ to $10^{-3}$), dropout rate (0.2-0.5), weight decay ($10^{-7}$ to $10^{-4}$), cluster embedding dimension (16, 32, 64), and number of clusters (3-7). The optimization objective was to maximize validation AUC.

Table~\ref{tab:optimization_results} compares the baseline and optimized model configurations. Optuna identified that increasing the number of LSTM layers from 3 to 4, doubling the batch size from 64 to 128, reducing dropout from 0.4 to 0.31, and slightly increasing the learning rate from $1.0 \times 10^{-4}$ to $2.25 \times 10^{-4}$ yielded substantial improvements. These changes increase model capacity while maintaining regularization through the optimized dropout rate. The larger batch size stabilizes gradients and enables the higher learning rate, accelerating convergence.

\begin{table}[h]
\centering
\caption{Hyperparameter Optimization Results}
\label{tab:optimization_results}
\begin{tabular}{@{}lrr@{}}
\toprule
\textbf{Hyperparameter} & \textbf{Baseline} & \textbf{Optimized} \\
\midrule
Validation AUC & 0.6947 & \textbf{0.7572} \\
LSTM Layers & 3 & 4 \\
Batch Size & 64 & 128 \\
Learning Rate & $1.0 \times 10^{-4}$ & $2.25 \times 10^{-4}$ \\
Dropout Rate & 0.4 & 0.311 \\
Weight Decay & $1.0 \times 10^{-5}$ & $7.56 \times 10^{-6}$ \\
Cluster Embed Dim & 32 & 32 \\
Number of Clusters & 5 & 5 \\
\midrule
\textbf{AUC Improvement} & & \textbf{+9.0\%} \\
\bottomrule
\end{tabular}
\end{table}

To validate the final optimized model on unseen data, I tested both the baseline and optimized models on 300 windows extracted from 100 confirmed exoplanet host stars from the TESS and Kepler missions. These are real planetary systems that were never used during training or validation, providing an independent test of generalization performance. As with the 7-star test, I processed the light curves through the standard pipeline (detrending, normalization, window extraction, BLS feature computation) and assigned each test window to its nearest pre-learned cluster center from the training checkpoint. No re-fitting of K-means was performed on this test data. Table~\ref{tab:planet_test_results} presents the results. The baseline model, despite achieving AUC 0.6947 on the validation set, made zero positive predictions (all 300 windows scored below 0.5), indicating poor calibration on real planet data. In contrast, the optimized model made 16 positive predictions (5.3 percent of windows) and achieved a maximum probability of 0.6666 for TIC 261337380, demonstrating improved decision boundary learning and calibration.

\begin{table}[h]
\centering
\caption{Testing on 100 Confirmed Exoplanet Systems (300 Windows)}
\label{tab:planet_test_results}
\begin{tabular}{@{}lrrrr@{}}
\toprule
\textbf{Model} & \textbf{Mean Prob} & \textbf{Max Prob} & \textbf{Positive Preds} & \textbf{Top Candidate} \\
\midrule
Baseline & 0.289 & 0.435 & 0/300 (0\%) & TIC 323517227 \\
Optimized & 0.236 & \textbf{0.667} & \textbf{16/300 (5.3\%)} & \textbf{TIC 261337380} \\
\bottomrule
\end{tabular}
\end{table}

The optimized model's top 10 candidates (Table~\ref{tab:top_candidates}) include several known exoplanet hosts with varying orbital periods and transit depths. For example, TIC 261337380 (max probability 0.667) and TIC 410417579 (max probability 0.644) both exceed the 0.5 decision threshold with high confidence. The model's ability to rank these confirmed planets highly despite never seeing them during training demonstrates that it has learned generalizable features of planetary transits rather than memorizing the training set. The variation in predictions across the three windows per star (shown in the "All 3 Windows" column) reflects the model's sensitivity to phase alignment and window positioning, with highest probabilities typically occurring for windows centered on transit events.

\begin{table}[h]
\centering
\caption{Top 10 Planet Candidates from Optimized Model}
\label{tab:top_candidates}
\begin{tabular}{@{}lrp{4.8cm}@{}}
\toprule
\textbf{TIC ID} & \textbf{Max Prob} & \textbf{All 3 Windows (phase 0.0, +0.25, -0.25)} \\
\midrule
261337380 & 0.667 & 0.433, \textbf{0.667}, 0.183 \\
410417579 & 0.644 & 0.187, 0.014, \textbf{0.644} \\
30415110 & 0.596 & 0.509, 0.479, \textbf{0.596} \\
323517227 & 0.587 & 0.226, 0.017, \textbf{0.587} \\
293525326 & 0.583 & 0.140, 0.007, \textbf{0.583} \\
375034775 & 0.576 & 0.247, 0.168, \textbf{0.576} \\
277911545 & 0.560 & 0.338, 0.356, \textbf{0.560} \\
409891588 & 0.548 & 0.427, \textbf{0.548}, 0.081 \\
128917881 & 0.544 & 0.174, \textbf{0.544}, 0.303 \\
100229772 & 0.527 & 0.273, 0.003, \textbf{0.527} \\
\bottomrule
\end{tabular}
\end{table}

\subsection{Transformer Approach}

\subsubsection*{Model Configuration}
The evaluated model employs a compact transformer architecture for use on a limited dataset (100 light curves) with the following specifications:
\begin{itemize}
    \item Embedding dimension: \(d_{model} = 64\)
    \item Number of encoder layers: \(n_{layers} = 2\)
    \item Number of attention heads: \(n_{heads} = 4\)
    \item Best validation loss: \(0.6263\)
\end{itemize}

\subsubsection*{Performance Metrics}

The model was evaluated on a held-out test set of 15 samples, achieving the following metrics:

\begin{table}[h]
\centering
\caption{Overall Classification Metrics}
\begin{tabular}{@{}lr@{}}
\toprule
\textbf{Metric} & \textbf{Score} \\
\midrule
Accuracy        & 0.5333 \\
Precision       & 0.5385 \\
Recall          & 0.8750 \\
F1-Score        & 0.6667 \\
ROC-AUC         & 0.6071 \\
\bottomrule
\end{tabular}
\end{table}

\subsubsection*{Per-Class Performance}

The classification report reveals substantial performance asymmetry between the two classes:

\begin{table}[h]
\centering
\caption{Classification Report by Class}
\begin{tabular}{@{}lcccc@{}}
\toprule
\textbf{Class} & \textbf{Precision} & \textbf{Recall} & \textbf{F1-Score} & \textbf{Support} \\
\midrule
Non-Transit & 0.5000 & 0.1429 & 0.2222 & 7 \\
Transit     & 0.5385 & 0.8750 & 0.6667 & 8 \\
\midrule
Macro avg.  & 0.5192 & 0.5089 & 0.4444 & 15 \\
Weighted avg. & 0.5205 & 0.5333 & 0.4593 & 15 \\
\bottomrule
\end{tabular}
\end{table}

\begin{figure}
    \centering
    \includegraphics[width=1.0\linewidth]{confusion_matrix.png}
    \caption{Confusion matrix for initial trial of Transformer training}
    \label{fig:placeholder}
\end{figure}

\subsubsection*{Analysis}

The model demonstrates a strong class imbalance in predictive behavior:
\begin{itemize}
    \item{High recall for Transit class (87.5\%):} The model successfully identifies most exoplanet transit events, minimizing false negatives.
    \item{Poor recall for Non-Transit class (14.3\%):} The model exhibits a tendency to over-predict the Transit class, resulting in numerous false positives.
    \item{Moderate overall accuracy (53.3\%):} Performance is only marginally better than random guessing for binary classification.
    \item{ROC-AUC of 0.607:} Indicates weak discriminative ability between classes.
\end{itemize}

The low F1-score for the Non-Transit class (0.22) suggests the model struggles to correctly classify negative samples, likely due to insufficient model capacity (\(d_{model}=64\), \(n_{layers}=2\)) or suboptimal hyperparameters. The small test set size (\(n=15\)) also limits the statistical reliability of these metrics.

\subsubsection*{Future Improvements}

To improve model performance:
\begin{enumerate}
    \item Increase model capacity by expanding \(d_{model}\) and \(n_{layers}\).
    \item Apply class-weighted loss functions to address prediction bias toward the Transit class.
    \item Extend hyperparameter tuning to explore larger learning rates and regularization parameters.
    \item Increase test set size for more robust evaluation statistics.
\end{enumerate}


\section*{Experimental Plan \& Milestones}
Week 1: Assemble the small curated subset (\textasciitilde{}10 positive, matched negatives) and write a minimal data-prep plan.

Week 2: Train a simple baseline (the small curated dataset) to set a floor; implement the first Machine Learning run.

Week 3: Tune individual MLs and pick the better or different variant based on validation F1/AUC.

Week 4: Add optional pseudo/synthetic injections to balance classes; re‑evaluate and do a brief error analysis.

Week 5–6: Polish: gather findings and start final metrics/tables/plots, and input to the write‑up for the team paper.

\section*{Overall Project Risks \& Mitigations}
• Too few positives: mitigate with pseudo/synthetic injections and careful cross‑validation.
• Overfitting the small set: use validation splits, early stopping, and simple models first.
• Data cleaning surprises: keep pre-processing minimal and documented; track all changes in the notebook.

\section*{Conclusion}
In conclusion, through this research, we plan to solve two problems: (1) Automatically classifying planets vs. non-planets in light curves. (2) Building models that are more generalizable to multiple datasets to reduce the need to create a model for each dataset released. The plan of action is for each member to develop a different machine learning model – CNN, RNN, and Transformer, which is trained on Kepler, TESS, and simulated data. We plan to compare these models using shared evaluation metrics, then analyze which approach is the most adaptable for future exoplanet discovery.

\section*{AI Use Disclosure}

Portions of this work were completed with assistance from \textbf{Claude Code (Sonnet 4.5)}, version \texttt{claude-sonnet-4-5-20250929}.

AI assistance for the RNN component included:
\begin{itemize}
\item Understanding RNN concepts from lecture materials and research papers
\item Implementing BiLSTM architecture with PyTorch
\item Debugging data pipeline and training code (handling Windows multiprocessing issues, GPU memory management)
\item Writing comprehensive documentation and comments in code
\item Experiment design and hyperparameter tuning guidance
\item LaTeX formatting and table generation
\end{itemize}

All code was reviewed, understood, and tested by the students before submission. All experimental results were obtained by running code on real data. All design decisions (architecture choices, hyperparameters, evaluation metrics) were made by the students. The AI did not complete the project autonomously -- student understanding, testing, debugging, and decision-making were central to all work.


\bibliography{resourceFile}

\end{document}
